\section{Replication Setup and Protocols}
\label{sec:setup}

\subsection{Datasets}

We use the same ten datasets as Farid et al.'s Table 5: breast cancer,
contact lenses, diabetes, glass, iris, soybean, vote, image segmentation,
tic-tac-toe, and NSL-KDD~\cite{tavallaee2009detailed}. NSL-KDD is the
corrected, de-duplicated successor to KDD Cup 1999 and is the one dataset
in the set built specifically to study intrusion-detection classifiers
rather than as a generic UCI benchmark; it is also the dataset on which a
preprocessing-leakage effect has previously been
documented~\cite{bouke2023empirical}, which is one reason its results are
reported separately rather than folded into an average.

\subsection{Two evaluation protocols}

Every experiment in this paper compares the same two protocols, applied to
whichever stages an arm uses (the NB filter of Algorithm~1, the tree-based
attribute weighting of Algorithm~2, or both):
\begin{itemize}
\item \textbf{Protocol B (refit-per-fold, honest).} For each of the 10
cross-validation folds, every stage (the filter, the attribute
selector, and the final classifier) is fit on that fold's training
split only, and scored on that fold's held-out test split. No stage ever
sees a label from a test instance before that instance is scored.
\item \textbf{Protocol A (fit-before-CV).} The filter and/or the attribute
selector are fit once on the entire dataset, using every instance's label.
Cross-validation is then run only over the final classifier, on the data
already cleaned or reduced by that one full-data fit. This is the protocol
consistent with a plain reading of Farid et al.'s Section 4, which never
mentions folds when it describes Algorithms 1 and 2, and only introduces
10-fold cross-validation afterward, in Section 5, as the way the resulting
classifier is scored.
\end{itemize}
Where a fully leaky upper bound is useful for calibration, we also report
\textbf{protocol C}: the stages and the final classifier are fit on the
whole dataset and then scored on that same data, with no held-out split at
all. Protocol C is not a protocol anyone would defend; it exists only to
show how high a number leakage can reach.

\subsection{R1: faithful replication with the paper's own components}

R1 reproduces Farid et al.'s Tables 8--11 as closely as the paper's own
tools allow. C4.5 is Weka J48 3.8.6~\cite{witten2016weka}, pruned, with
Weka's default options (\texttt{-C 0.25 -M 2}), matching Quinlan's original
C4.5~\cite{quinlan1993c45} far more closely than a CART-style tree does,
because J48 splits a nominal attribute into one branch per value rather
than into binary one-hot tests; this matters because Algorithm~2's notion
of ``the attribute'' and its minimum-depth weight are only meaningful in
the attribute space the paper used. NB is a from-scratch implementation
(add-one Laplace counts on nominal attributes, Gaussian likelihood on
numeric ones) fit on the datasets' original, non-one-hot columns, because
Weka's own NB diverges from this textbook form on attributes with very
wide numeric ranges such as NSL-KDD's byte counts. Algorithm~1 uses this
NB as its noise judge; Algorithm~2 uses J48's tree structure for the
minimum-depth weights and the same NB, weighted by Eq.~\eqref{eq:14}, as
its final classifier. Both protocols above are run for every arm, with 3
seeds and stratified 10-fold cross-validation per seed. As a sensitivity
check, R1 is repeated once with J48 unpruned (\texttt{-U -M 2}).

\subsection{E3a: an independent leakage audit}

E3a re-runs the same protocol-A/B/C comparison with a second, independent
implementation: scikit-learn~\cite{pedregosa2011scikit} decision trees
grown with entropy splitting (CART, not C4.5), a one-hot encoding of
nominal attributes, and a NB with mixed likelihoods (Bernoulli on the
0/1 one-hot columns, Gaussian elsewhere). If the protocol-A inflation seen
in R1 were an artifact of the Weka/J48 implementation, it should not
reappear here; if it is a property of the algorithms and the leaky
protocol, it should. E3a uses 3 seeds and stratified 10-fold
cross-validation, tree pruning strength $\mathrm{ccp\_alpha}=0.01$.

\subsection{E9: the two orders of Algorithm 1 and Algorithm 2}

E9 chains Algorithm~1 (instance filtering, written $N$) and Algorithm~2
(attribute weighting, written $A$) in both possible orders, always under
protocol B: Path~1 is $N\to A$ (filter, then weight attributes on the
filtered rows), Path~2 is $A\to N$ (weight attributes, then filter
instances using the weighted NB as judge). Both paths end in either a
final NB or a final decision tree, alongside single-stage (Algorithm~1
alone, Algorithm~2 alone) and no-cleaning baseline arms, all evaluated on
the same folds so every arm is directly comparable. E9 uses the same
scikit-learn CART/one-hot/mixed-NB implementation as E3a, tree pruning
strength $\mathrm{ccp\_alpha}=0$ (no pruning), and Algorithm~2's weights
are also handed to Algorithm~1's judge and to the final NB when both
stages run. E9 uses 5 seeds and stratified 10-fold cross-validation.

\subsection{Statistical testing}

Every comparison in this paper is a Wilcoxon signed-rank test over the ten
datasets: for a given pair of conditions, the ten paired accuracy
differences (one per dataset, averaged over that condition's seeds and
folds already) are tested against a null median difference of zero,
two-sided, using SciPy. With only ten datasets, this gives $n=10$ paired
observations per test, which is a small sample by any standard and a low-power
test in the sense discussed by Dem\v{s}ar~\cite{demsar2006statistical};
we report the Wilcoxon statistic, the raw $p$-value, the median paired
difference, and a Holm step-down correction applied across the family of
eight tests reported in Table~\ref{tab:tests}, and we do not claim
significance where the corrected $p$-value does not support it. No
per-fold results are stored in the underlying CSV files, so no
resampled or corrected paired $t$-test is computed; only the
between-dataset Wilcoxon test is reported.
